\begin{abstract}
Before a batch pod runs on Kubernetes, someone has to decide how much memory and CPU it gets and whether
it gets a GPU. Too little makes it fail or run slowly, and too much wastes the cluster. We present
RightSize Bench, a benchmark that scores this decision in dollars: each request is replayed against a
measured run, failed attempts and throttling included, and compared with the cheapest request that would
have run it. The benchmark has 32,082 measured runs of 33 families of generated Python tasks that use real
libraries, from dataframes and SQL to PyTorch training and sharded language models. Every run appears
twice, as a raw Kubernetes Job whose inputs are staged in one of four ways, and as a task submitted
through Flyte, an orchestrator whose runtime we measure at about 70\,MiB per pod. We checked the Job labels
in about 480 real pods, which found two errors in our first release.

We release the benchmark with a set of baseline models: gradient-boosted quantile trees, and a 1.7B
decision model whose pointer heads score the request options written in its prompt, trained with
cross-entropy and then with reinforcement learning against the cost. On new tasks from familiar families,
the trees pooled with the decision model come within \rTestPoolDec\% of the cheapest request on Job pods,
against \rTestSized\% for a rule that scales earlier runs by input size and reads the pod spec, and
\rTestZS\% for a zero-shot 27B model. On seven held-out families every approach does much worse, and the
rule (\rUnseenSized\%) is as good as anything we trained. Models trained on Jobs move their memory request
with the volume medium in the manifest about three times in four. \codeSeenShare\% of the test tasks share
their exact code with a training task, so the held-out families are the benchmark's test of new code.
\end{abstract}
